
\documentclass[11pt]{article}

\usepackage[a4paper,margin=0.9in]{geometry}
\usepackage[T1]{fontenc}
\usepackage[utf8]{inputenc}
\usepackage{setspace}
\usepackage{hyperref}
\usepackage{enumitem}
\usepackage{lmodern}
\usepackage{natbib}

\onehalfspacing

\title{Etomidate, ketamine, and propofol for rapid sequence intubation in critically ill adults:\\
Protocol for a systematic review and network meta-analysis of randomized trials}

\author{Fernando G.\ Zampieri \and Raysa C Schmidt \and Bruno AMP Besen \and Fernando JDS Ramos \and Fl\'avio GR Freitas \and Fl\'avia R.\ Machado\\[4pt]
\textnormal{for the PROMINE Investigators}}

\date{December 10, 2025}

\begin{document}

\maketitle

\section*{Registration}

This protocol will be registered in the International Prospective Register of Systematic Reviews (PROSPERO). The registration ID will be added when available.

\section{Background}

Emergency tracheal intubation in critically ill adults is a high-risk procedure in which the choice of induction agent may substantially affect hemodynamic stability and patient outcomes. Etomidate, ketamine, and propofol are widely used hypnotic agents for rapid sequence induction (RSI), and ketamine--propofol combinations (``ketofol'') are also used in some settings. These agents differ in their cardiovascular effects, and clinicians often select them based on perceived hemodynamic profiles rather than comparative effectiveness data.

Recent large randomized trials provide new, high-quality evidence that has not yet been synthesized within a unified framework. Existing systematic reviews have focused mainly on pairwise comparisons and have not integrated all relevant agents using network meta-analysis (NMA). Relevant new trials incluse \cite{PROMINE2025,RSI2025}.

A comprehensive synthesis of randomized evidence is needed to clarify the comparative safety and effectiveness of induction agents used for emergency tracheal intubation in critically ill adults, particularly with respect to hemodynamic instability and mortality.

\section{Objectives}

The objectives of this review are:

\begin{enumerate}[label=\arabic*.]
  \item To compare the effects of etomidate, ketamine, propofol, and ketofol on outcomes of emergency tracheal intubation in critically ill adults.
  \item To determine the relative impact of these agents on short-term mortality, post-induction hypotension, cardiovascular collapse, vasopressor initiation, and first-pass intubation success.
  \item To use network meta-analysis to estimate comparative effectiveness across all eligible induction agents using both direct and indirect randomized evidence.
  \item To incorporate new randomized data (PROMINE and RSI) to provide an up-to-date assessment of induction agent performance.
  \item To explore sources of heterogeneity, including clinical setting (emergency department, intensive care unit, prehospital), baseline hemodynamics, and co-interventions, if possible.
\end{enumerate}

\section{Methods}

\subsection{Eligibility criteria}

\subsubsection{Population}

The following inclusion and exclusion criteria will be applied.

\paragraph{Inclusion}

\begin{itemize}
  \item Adults (typically $\geq$16 or $\geq$18 years, as defined by individual studies).
  \item Critically ill or acutely ill adults undergoing \emph{emergency} or \emph{rapid sequence} tracheal intubation.
  \item Intubation performed in emergency departments, intensive care units, acute hospital wards, or prehospital / EMS settings.
\end{itemize}

\paragraph{Exclusion}

\begin{itemize}
  \item Elective or planned operating-room intubations.
  \item Procedural sedation without tracheal intubation.
  \item Obstetric anesthesia (for example, cesarean section induction).
  \item Pediatric studies in which adult data cannot be isolated.
  \item Studies focused on sedation after intubation rather than induction.
  \item Elective surgical populations without acute critical illness.
\end{itemize}

\subsubsection{Interventions and comparators}

\paragraph{Eligible induction agents}

\begin{itemize}
  \item Etomidate.
  \item Ketamine.
  \item Propofol.
  \item Ketofol (ketamine--propofol combination).
\end{itemize}

Trials must randomize the hypnotic agent administered immediately prior to laryngoscopy for emergency or rapid sequence tracheal intubation. Any eligible agent may serve as comparator. Co-interventions (for example, opioids, neuromuscular blockers, preoxygenation strategies) are allowed if applied similarly across arms.

\paragraph{Excluded interventions}

\begin{itemize}
  \item Trials in which the randomized intervention is not the hypnotic induction agent (for example, neuromuscular blocker, airway device, fluid strategy).
  \item Sedative agents used solely for post-intubation ICU sedation.
  \item Agents not used in emergency RSI in critically ill adults (for example, remimazolam or ciprofol in elective anesthesia).
\end{itemize}

\subsubsection{Outcomes}

Trials must report at least one of the following outcomes or provide sufficient data to derive it.

\paragraph{Primary outcome}

\begin{itemize}
  \item \textbf{Short-term mortality}: in-hospital or 28--30 day mortality (the measure closest to 28--30 days will be used). The effect measure will be the odds ratio (OR) with 95\% confidence interval (CI).
\end{itemize}

\paragraph{Secondary outcomes}

\begin{itemize}
  \item \textbf{Cardiovascular collapse}: composite including hypotension, new or increased vasopressor requirement, or peri-intubation cardiac arrest, as defined by individual studies. Where reported, individual components will be extracted and analyzed separately. Effect measure: OR with 95\% CI.
  \item \textbf{Post-induction hypotension}: blood pressure below study-defined thresholds or a major decline from baseline within 0--15 minutes after induction. Effect measure: OR with 95\% CI.
  \item \textbf{Vasopressor initiation or escalation}: new or increased vasopressor requirement. Data will be extracted at multiple time points where available (up to 30 minutes, 1 hour, and 24 hours after induction). Effect measure: OR with 95\% CI.
  \item \textbf{First-pass intubation success}: successful endotracheal tube placement on the first laryngoscopy attempt. Effect measure: OR with 95\% CI.
\end{itemize}

Outcomes reported at different time points will be analyzed separately and not pooled across time windows.

No other outcomes will be analyzed.

\subsubsection{Study designs}

\paragraph{Inclusion}

\begin{itemize}
  \item Randomized controlled trials (parallel-group).
\end{itemize}

\paragraph{Exclusion}

\begin{itemize}
  \item Observational studies (cohort, case-control, cross-sectional).
  \item Case series or case reports.
  \item Cross-over trials.
  \item Simulation or volunteer studies.
  \item Conference abstracts without sufficient data for risk of bias assessment and effect estimation.
\end{itemize}

\subsection{Information sources}

The following databases will be searched from inception:

\begin{itemize}
  \item MEDLINE (via PubMed)
  \item Embase (through Ovid)
\end{itemize}

Additional identification:

\begin{itemize}
  \item Forward citation searching (“snowballing”) of included trials
  \item Screening trial registries (ClinicalTrials.gov and WHO ICTRP) for contextual information on completed or ongoing trials
\end{itemize}

Only randomized trials will be included. Only studies published in English will be considered.

\subsection{Search strategy}

A draft MEDLINE search strategy is as follows and will be refined with information specialist input:

\begin{verbatim}
(
  "Intubation, Intratracheal"[Mesh] OR intubat*[tiab] OR "tracheal intubation"[tiab] OR
  "rapid sequence"[tiab] OR "rapid-sequence"[tiab] OR RSI[tiab]
)
AND
(
  "Emergency Service, Hospital"[Mesh] OR "Intensive Care Units"[Mesh] OR
  emergency[tiab] OR "emergency department"[tiab] OR ED[tiab] OR
  "critical illness"[Mesh] OR "critically ill"[tiab] OR ICU[tiab] OR
  "prehospital"[tiab]
)
AND
(
  etomidate[tiab] OR ketamine[tiab] OR propofol[tiab] OR
  amidate[tiab] OR diprivan[tiab] OR ketofol[tiab]
)
AND
(
  randomized controlled trial[pt] OR controlled clinical trial[pt] OR
  random*[tiab] OR trial[tiab] OR "clinical trial"[pt]
)
NOT
(
  animals[mh] NOT humans[mh]
)
\end{verbatim}

The initial Embase search is:

\begin{verbatim}
1. exp endotracheal intubation/ or exp rapid sequence induction/
2. (intubat* or "tracheal intubation" or "rapid sequence" or "rapid-sequence" or RSI).ti,ab.
3. 1 or 2
4. exp emergency ward/ or exp intensive care unit/ or exp critical illness/
5. (emergency or "emergency department" or ED or "critically ill" or ICU or prehospital).ti,ab.
6. 4 or 5
7. (etomidate or ketamine or propofol or amidate or diprivan or ketofol).ti,ab.
8. exp etomidate/ or exp ketamine/ or exp propofol/
9. 7 or 8
10. 3 and 6 and 9
11. randomized controlled trial/ or controlled clinical trial/
12. (random* or trial or RCT).ti,ab.
13. 11 or 12
14. 10 and 13
15. (animal/ or nonhuman/) not human/
16. 14 not 15
\end{verbatim}

Final search strategies for all databases will be reported in an appendix to the main manuscript.

\subsection{Study selection}

Search results will be imported into reference management software and then into a systematic review platform for de-duplication and screening. Two reviewers (or a person/machine combination with human verification) will independently screen titles and abstracts for potential eligibility. Full texts of potentially eligible reports will then be assessed independently using the predefined criteria.

Disagreements at any stage will be resolved through discussion or, if necessary, by consultation with a third reviewer. Reasons for exclusion at the full-text stage will be recorded. The study selection process will be summarized in a PRISMA flow diagram.

\subsection{Data extraction}

Two reviewers (or a person/machine combination with human verification) will independently extract data using a standardized, piloted form. Extracted information will include:

\begin{itemize}
  \item Study characteristics: authors, year, country, setting (ED, ICU, prehospital), single versus multicenter, funding source.
  \item Population: inclusion and exclusion criteria, sample size, age, sex, primary diagnosis, baseline severity scores, baseline hemodynamics and vasopressor use.
  \item Interventions and comparators: induction agent, dose, timing, co-induction or adjunct agents, neuromuscular blocker type and dose, preoxygenation strategy, and RSI protocolization.
  \item Outcomes: definitions, time windows, and arm-level results for all eligible endpoints.
  \item Risk of bias domains (see below).
\end{itemize}

No contact with study authors is planned. If multiple reports refer to the same trial, they will be collated and treated as a single study, using the most complete and recent data.

\subsection{Risk of bias assessment}

Two reviewers (with optional machine assistance and full human verification) will independently assess risk of bias using the Cochrane Risk of Bias 2 (RoB~2) tool, considering:

\begin{itemize}
  \item Bias arising from the randomization process.
  \item Bias due to deviations from intended interventions.
  \item Bias due to missing outcome data.
  \item Bias in measurement of the outcome.
  \item Bias in selection of the reported result.
\end{itemize}

Judgements (low risk, some concerns, high risk) will be incorporated into sensitivity and subgroup analyses.

\subsection{Effect measures}

For binary outcomes (mortality, hypotension, cardiovascular collapse, vasopressor initiation, first-pass success), odds ratios with 95\% CIs will be calculated from arm-level data. Where necessary, alternative effect measures will be converted to ORs using standard methods.

Continuous outcomes, if analyzed descriptively (for example, minimum mean arterial pressure), will be summarized using mean differences or standardized mean differences with 95\% CIs. Continuous outcomes will not be the focus of the NMA.

\subsection{Data synthesis}

\subsubsection{Pairwise meta-analysis}

Where there are at least two sufficiently similar trials, conventional pairwise meta-analyses will be performed using random-effects models. Statistical heterogeneity will be quantified using $\tau^2$ and $I^2$, and 95\% prediction intervals will be presented where appropriate.

\subsubsection{Network meta-analysis}

A random-effects NMA in a frequentist framework (for example, using the \texttt{netmeta} package in R) will be used to estimate relative treatment effects between all pairs of eligible induction agents for each outcome.

The NMA will assume transitivity; this assumption will be evaluated by comparing potential effect modifiers---such as setting, baseline hemodynamics, and co-interventions---across treatment comparisons. Consistency between direct and indirect evidence will be assessed using:

\begin{itemize}
  \item A global design-by-treatment interaction model.
  \item Local node-splitting or analogous approaches to compare direct and indirect estimates.
\end{itemize}

Treatment ranking will be summarized using P-scores or surface under the cumulative ranking (SUCRA) values, with appropriate caution regarding uncertainty and evidence quality.

\subsection{Heterogeneity and inconsistency}

Heterogeneity in pairwise meta-analyses will be assessed using $\tau^2$ and $I^2$. In NMA, the between-study variance for each outcome will be estimated and explored through subgroup and sensitivity analyses when data permit.

Predefined potential effect modifiers include:

\begin{itemize}
  \item Clinical setting (ED versus ICU versus prehospital).
  \item Baseline vasopressor use.
  \item Degree of RSI protocolization (strict protocol versus pragmatic practice).
  \item Overall risk of bias.
\end{itemize}

\subsection{Subgroup and sensitivity analyses}

Planned subgroup analyses include:

\begin{itemize}
  \item Setting (ED, ICU, prehospital).
  \item Presence versus absence of vasopressor use at baseline.
  \item High versus low risk of bias trials.
\end{itemize}

Sensitivity analyses will consider:

\begin{itemize}
  \item Excluding high risk of bias trials.
  \item Excluding small trials (for example, fewer than 30 participants per arm).
  \item Collapsing or removing ketofol as a separate node to assess its impact on network geometry and estimates.
\end{itemize}

\subsection{Reporting bias assessment}

Risk of bias due to missing results will be assessed by evaluating selective non-reporting of outcomes and, when there are at least ten studies for an outcome, by inspecting funnel plots and considering statistical tests for small-study effects. No author contact is planned to obtain missing data.

\subsection{Certainty of evidence}

The certainty of the evidence for key comparisons and outcomes will be assessed using the GRADE framework adapted for NMA (for example, using CINeMA), taking into account risk of bias, inconsistency, indirectness, imprecision, and publication bias. Summary-of-findings tables will present results for the most clinically relevant comparisons (for example, etomidate versus ketamine, ketamine versus propofol, etomidate versus propofol).

\subsection{Ethics and dissemination}

As this review uses only published aggregate data, formal ethics approval is not required. Findings will be disseminated through peer-reviewed publications and scientific meetings. Data and analysis code will be made available in an appropriate public repository, subject to journal policies.

\section*{Funding and conflicts of interest}

The review has no external funding. Institutional support is provided by participating authors' institutions. Any conflicts of interest will be fully disclosed in the final manuscript.

\bibliographystyle{plainnat} % or abbrvnat, or whatever you like
\bibliography{references}
\end{document}
